\input{template.tex}

\ifPDFTeX
  \renewcommand{\rmdefault}{cmr}
\fi

\def\SheetNo{9}
\setlength{\emergencystretch}{2em}

\begin{document}

\sheettitle{Связность}

\begin{problems}

\item Для простого графа $G$ на $n\ge3$ вершинах выписали все $n$ графов $G-v$, полученных удалением по одной вершине. Имена вершин стёрли, одинаковые графы оставили с учётом кратности. Докажите, что по этому списку можно восстановить весь исходный граф с точностью до переименования вершин.

\item Докажите, что вершинам любого дерева с $m$ рёбрами можно приписать числа $0,1,\ldots,m$ по одному разу так, чтобы модули разностей чисел на концах рёбер были попарно различны.

\item Зафиксируем простой граф $H$. Пусть в графе $G$ на $n\ge1$ вершинах никакое множество вершин со всеми рёбрами между ними не образует граф, изоморфный $H$. Докажите, что для некоторого $c>0$, зависящего только от $H$, в любом таком графе $G$ найдётся хотя бы $n^c$ попарно смежных или хотя бы $n^c$ попарно несмежных вершин.

\item Докажите, что рёбра любого связного графа на $n$ вершинах можно покрыть не более чем $\lceil n/2\rceil$ простыми путями так, чтобы каждое ребро вошло ровно в один путь.

\item Пусть $c(G-S)$ --- число компонент, оставшихся после удаления множества вершин $S$. Найдётся ли число $t\ge1$ со следующим свойством? Если в связном графе $G$ на хотя бы трёх вершинах для каждого $S$ с $c(G-S)>1$ выполняется $|S|\ge t\,c(G-S)$, то в $G$ есть гамильтонов цикл.

\item В связном графе выбрали три простых пути максимальной длины. Докажите, что у них есть общая вершина.

\item В ориентированном графе на $n\ge1$ вершинах нет петель, кратных и встречных дуг. Из каждой вершины выходит не менее $n/3$ дуг. Докажите, что в графе есть направленный цикл длины $3$.

\item В связном графе любую вершину $u$ можно перевести в любую вершину $v$ перенумерацией, сохраняющей смежность. Докажите, что в таком графе есть гамильтонов путь.

\item Пусть $k\ge2$ и в графе больше $k$ вершин. Удаление любого множества из менее чем $k$ вершин сохраняет его связность. Докажите, что любые два цикла максимальной длины в этом графе имеют не менее $k$ общих вершин.

\item Удаление любого одного ребра связного графа сохраняет его связность. Докажите, что рёбрам графа можно назначить направления и веса из множества $\{1,2,3,4\}$ так, чтобы у каждой вершины сумма весов входящих рёбер равнялась сумме весов исходящих.

\end{problems}

\end{document}
